\begin{textsample}{POS Dim 3 – 2023\_c5 – Score 4.00 – 2023\_c5/enade\_2023\_eng\_prod\_c677c7f341.txt}  \label{ex:f3_pos_060}
CAPÍTULO 4 DISTRIBUIÇÃO DE CONCEITOS A finalidade deste capítulo é apresentar um panorama nacional sobre a distribuição do Conceito Enade 2023.
Para \textbf{traçar} esse panorama, recorre -se a tabelas e gráficos, que são acompanhados de comentários e análises que contribuem para a contextualização dos dados neles expostos.
Para visualizar o panorama do Enade 2023, em nível nacional, na primeira seção deste capítulo, os dados sobre o número de cursos participantes são organizados por grande região.
Na segunda, o mapeamento do conceito do exame é feito por categoria administrativa das IES e modalidade de oferta dos cursos, de maneira cruzada às grandes regiões.
Na última seção, a lógica de apresentação dos dados é feita por organização acadêmica e, também, por grande região de funcionamento dos cursos, no país como um todo. 4.1 PANORAMA NACIONAL DA DISTRIBUIÇÃO DOS CONCEITOS A lógica adotada na apresentação dos dados, nesta seção e nas duas seguintes, pressupõe o registro de duas informações de natureza reiterativa.
O primeiro diz respeito aos valores possíveis do Enade, que variam em níveis de 1 a 5, sendo que, à medida que esses valores aumentam, o curso é melhor avaliado, e ao fato de que aqueles cursos que tiveram somente um ou nenhum concluinte participante não obtêm o Conceito Enade, ficando sem conceito ( SC ), como ressaltado na apresentação deste relatório.
O segundo registro é que, caso haja diferença entre os cursos tabulados neste capítulo e no capítulo 1, ela corresponderá aos cursos que não foram avaliados, ou seja, àqueles, em princípio, sem estudantes concluintes que tivessem sido inscritos para a prova.
Na Tabela 4.1, são apresentadas a quantidade e a distribuição de cursos de Engenharia de Produção participantes do Enade 2023, por faixa de conceito e grande região.
A partir dos dados da Tabela 4.1, nota -se que, dos 638 cursos participantes, 239 ( 37,5 \% ) classificaram -se com conceito 3, o valor modal.
Esse foi também o conceito modal para as regiões Sul e Sudeste.
Por região, o percentual de cursos que se classificaram com conceito 3 foram : Centro-Oeste ( 29,4 \% ), Nordeste ( 28,1 \% ), Norte ( 22,2 \% ), Sudeste ( 40,1 \% ) e Sul ( 42,4 \% ).
O conceito 2 foi o segundo mais frequente em nível nacional ( 31,0 \%, correspondendo a 198 cursos ), tendo sido o segundo conceito mais frequente nas regiões Sul e Sudeste.
O conceito 4 foi o terceiro ( 16,1 \%, correspondendo a 103 cursos ).
Houve, ainda, 27 ( 4,2 \% ) cursos que receberam o conceito 5 e oito ( 1,3 \% ) cursos que receberam o conceito 1.
Dos 638 cursos de Engenharia de Produção, 63 ( 9,9 \% ) ficaram sem conceito ( SC ).
A região Norte participou com 27 cursos, o que corresponde a 4,2 \% do total nacional.
Desses, dez cursos receberam o conceito 2, o que equivale a 37,0 \% do total regional, o conceito modal na região.
A seis cursos ( 22,2 \% ) atribuiu -se o conceito 3 ; a três cursos ( 11,1 \% ), o conceito 4 ; e a dois cursos ( 7,4 \% ) o conceito 5.
Ainda nessa região, um curso ( 3,7 \% ) obteve conceito 1 e cinco cursos ( 18,5 \% ) ficaram sem conceito ( SC ).
A região Nordeste participou com 96 cursos, correspondentes a 15,0 \% do total nacional.
Nessa região, 39 cursos ( 40,6 \% em termos regionais ) obtiveram conceito 2, o conceito modal para a região.
Os conceitos 3 e 4 foram atribuídos, respectivamente, a 27 e a 15 cursos ( 28,1 \% e 15,6 \%, respectivamente ).
Ainda nessa região, três cursos ( 3,1 \% ) foram avaliados com o conceito 5 e dois cursos ( 2,1 \% ) foram avaliados com o conceito 1.
Outros dez cursos ( 10,4 \% ) ficaram sem conceito ( SC ).
Para os 337 cursos participantes da região Sudeste ( 52,8 \% do total nacional ), o conceito 3 foi o conceito modal, como já anteriormente assinalado, tendo sido obtido por 135 cursos ( 40,1 \% ).
O conceito 2 foi atribuído a 101 cursos ( 30,0 \% ) e o conceito 4, a 51 cursos ( 15,1 \% ).
Receberam os conceitos 5 e 1, respectivamente, 16 cursos ( 4,7 \% ) e três cursos ( 0,9 \% ).
Nessa região, 31 cursos ( 9,2 \% ) ficaram sem conceito ( SC ).
Os 144 cursos da região Sul corresponderam a 22,6 \% do total nacional.
Como já indicado, houve predominância do conceito 3 ( conceito modal ), atribuído a 42,4 \% dos cursos da região, ou seja, atribuído a 61 dos 144 cursos participantes na região Sul.
O conceito 2 foi atribuído a 35 cursos ( 24,3 \% ) ; o conceito 4, 30 cursos ( 20,8 \% ) ; e o conceito 5, a cinco cursos ( 3,5 \% ).
Na região Sul, um curso ( 0,7 \% ) obteve conceito 1 e 12 cursos ( 8,3 \% ) ficaram sem conceito ( SC ).
Já dos 34 cursos participantes na região Centro-Oeste ( 5,3 \% do total nacional ), 13 ( 38,2 \% em termos regionais ) receberam conceito 2, o conceito modal, dez cursos ( 29,4 \% ) obtiveram o conceito 3, quatro cursos ( 11,8 \% ) obtiveram o conceito 4 e um curso ( 2,9 \% ), o conceito 5.
Nessa região, um curso ( 2,9 \% ) obteve o conceito 1 e cinco cursos ( 14,7 \% ) ficaram sem conceito ( SC ).
Em um gráfico de distribuição cumulativa, em cada valor de abscissa ( eixo x ), o valor apresentado no eixo das ordenadas ( eixo y ) é o valor acumulado da variável ( Conceito Enade ).
Em linhas gerais, em um gráfico desta mesma tipologia, poligonais mais à esquerda representam, no caso dos dados apresentados nesta seção, as grandes regiões com a pior distribuição, e poligonais mais à direita, as grandes regiões com os melhores conceitos.
A partir destas diretrizes, no Gráfico 4.1, é apresentada a distribuição cumulativa do Conceito Enade, segundo a grande região.
No Gráfico 4.1, observando a linha vertical pontilhada, que \textbf{ajuda} a visualizar o percentual de conceitos acumulados até o valor 3, ou seja, conceitos 1, 2 e 3, pode -se notar, por exemplo, que a região Centro-Oeste ( linha \textbf{laranja} ) apresenta 82,8 \% dos cursos com conceito 3 ou abaixo deste valor.
A região Sul ( linha roxa ), por outro lado, apresenta 73,5 \% dos cursos na mesma situação : com conceito 3 ou abaixo.
A região Centro-Oeste apresenta os cursos com conceitos mais baixos, ao passo que a região Sul apresenta a distribuição com valores maiores.
Os cursos em IES situadas na região Sudeste ( linha vermelha ) apresentam uma situação intermediária. 4.2 CONCEITOS POR CATEGORIA ADMINISTRATIVA, MODALIDADE DE OFERTA DOS CURSOS E GRANDE REGIÃO Na Tabela 4.2, é apresentada a distribuição dos cursos de Engenharia de Produção participantes do Enade 2023, por categoria administrativa e modalidade de oferta, de acordo com os conceitos por eles alcançados, segundo as grandes regiões brasileiras.
De acordo com as informações da Tabela 4.2, dos 638 cursos participantes, 121 ( 19,0 \% ) eram ministrados em instituições públicas e 517 ( 81,0 \% ) em instituições privadas.
Quanto à modalidade de oferta, 573 ( 89,8 \% ) cursos eram presenciais e 65 ( 10,2 \% ), a distância.
De acordo com as informações da mesma tabela, em termos nacionais, entre instituições públicas e privadas, observa -se que, dos 27 cursos avaliados com conceito 5, 24 foram oferecidos em IES públicas, e 3, em IES privadas.
Dos 517 cursos participantes de IES privadas, o conceito 3 foi o valor modal, atribuído a 200 ( 38,7 \% ) cursos.
Em relação aos demais cursos de instituições privadas, 193 cursos ( 37,3 \% ) foram avaliados com o conceito 2, 52 cursos ( 10,1 \% ) foram avaliados com o conceito 4, oito cursos ( 1,5 \% ) foram avaliados com o conceito 1 e, três cursos ( 0,6 \% ), com o conceito 5.
Nessa categoria, 61 cursos ( 11,8 \% ) ficaram sem conceito ( SC ).
Na rede pública, o conceito modal foi o 4, atribuído a 51 cursos, correspondentes a 42,1 \% dos 121 cursos da categoria.
Dos demais cursos participantes, 39 cursos ( 32,2 \% ) obtiveram o conceito 3, 24 cursos ( 19,8 \% ) receberam o conceito 5 e cinco cursos ( 4,1 \% ), o conceito 2.
Nenhum curso ( 0,0 \% ) obteve o conceito 1.
Nessa categoria administrativa, dois cursos ( 1,7 \% ) ficaram sem conceito ( SC ).
Dos 573 cursos participantes da modalidade presencial, o conceito 3 foi o valor modal, atribuído a 217 ( 37,9 \% ) cursos.
Dos demais cursos dessa modalidade de oferta, 164 cursos ( 28,6 \% ) receberam o conceito 2 ; 101 cursos ( 17,6 \% ) receberam o conceito 4 ; 27 cursos ( 4,7 \% ), o conceito 5 ; e sete cursos ( 1,2 \% ), o conceito 1.
Nessa modalidade de oferta, 57 cursos ( 9,9 \% ) ficaram sem conceito ( SC ).
Dos 65 cursos participantes da modalidade a distância, o conceito 2 foi o valor modal, atribuído a 34 ( 52,3 \% ) cursos.
Dos demais cursos ofertados nessa modalidade, 22 cursos ( 33,8 \% ) receberam conceito 3, dois cursos ( 3,1 \% ) receberam o conceito 4 ; um curso ( 1,5 \% ) recebeu o conceito 1 ; e nenhum curso ( 0,0 \% ) recebeu conceito 5.
Nessa modalidade de oferta, seis cursos ( 9,2 \% ) ficaram sem conceito ( SC ).
Na análise por região, observa -se, que, na região Norte, as instituições públicas participaram com seis cursos ( 22,2 \% do total regional ), dos quais nenhum ficou sem conceito ( SC ) e um curso recebeu o conceito 5.
O conceito modal foi 3, atribuído a três cursos, correspondendo a 50,0 \% dessa categoria na região.
O conceito 4 teve uma concentração muito próxima do modal, com dois cursos ( 33,3 \% ), e o conceito 5 teve um curso ( 16,7 \% ).
Nenhum curso recebeu os conceitos 2 e 1.
As instituições privadas participaram com 21 cursos na região Norte ( 77,8 \% do total regional ), e cinco cursos ficaram sem conceito ( SC ).
O conceito modal nas instituições privadas na região foi o 2, atribuído a dez cursos, o que corresponde a 47,6 \% dessa categoria na região.
Os conceitos 3, 4 e 5 foram atribuídos a três, um e um cursos, respectivamente.
O conceito 1 foi atribuído a um curso ( 4,8 \% ).
Dos cursos oferecidos na região Norte, 23 eram presenciais e quatro ficaram sem conceito ( SC ).
Dos 23 cursos oferecidos, sete obtiveram o conceito 2 ( modal ) ; seis, o conceito 3 ; três, o conceito 4 ; dois, o conceito 5 ; um, o conceito 1 e quatro cursos ficaram sem conceito ( SR ).
Quatro cursos eram na modalidade a distância,tendo três recebido o conceito modal 2 e um curso ficou sem conceito ( SC ).
Nenhum curso recebeu os demais conceitos.
Na região Nordeste, a rede privada concentra 71 dos 96 cursos participantes, o equivalente a 74,0 \% do total da região.
Entre os cursos oferecidos por IES privadas no Nordeste, 39 cursos foram avaliados com conceito 2, o conceito modal ; 17 cursos obtiveram o conceito 3 ; cinco cursos, o conceito 4 ; dois cursos, o conceito 1 ; e nenhum curso obteve o conceito 5.
Nessa categoria administrativa e grande região, oito cursos ficaram sem conceito ( SC ).
As instituições públicas da região Nordeste participaram com 25 cursos ( 26,0 \% do total da região ).
Desses, o valor bimodal foram os conceitos 4 e 3, ambos com dez ( 40,0 \% ) cursos.
O conceito 5 foi atribuído a três cursos ; o conceito 2, a nenhum curso ; e o conceito 1, também, a nenhum curso.
Dois cursos ficaram sem conceito ( SC ).
Dos cursos oferecidos na região Nordeste, todos ( 91 ) eram presenciais : dez ficaram sem conceito ( SC ), dois obtiveram o conceito 1 ; 37, o conceito 2, o valor modal ; 24, o conceito 3 ; 15, o conceito 4 ; e três, o conceito 5.
Cinco cursos eram na modalidade a distância, tendo três cursos recebido o conceito modal 3 ; e dois, o conceito 2.
Nenhum curso ficou sem conceito ( SC ) ou recebeu os demais conceitos na região Nordeste e na modalidade a distância.
Na região Sudeste, a proporção de cursos da rede privada, 84,0 \%, correspondeu a 283 dos 337 cursos participantes.
Nessa combinação de categoria administrativa e grande região, o conceito modal foi o 3, atribuído a 120 cursos ; e 31 cursos ficaram sem conceito ( SC ).
Os demais foram avaliados com os seguintes conceitos : 2 ( 100 cursos ), 4 ( 27 cursos ), 1 ( três cursos ) e 5 ( dois cursos ).
Dos 54 cursos oferecidos por instituições públicas na região Sudeste, 14 obtiveram o conceito 5 ; 24, o conceito 4 ; 15, o conceito 3 ; um obteve conceito 2 e nenhum obteve o conceito 1, assim como nenhum curso ficou sem conceito ( SC ).
Ainda na região Sudeste, grande parcela dos cursos ( 92,6 \% ) foi oferecida na modalidade presencial : 312 dos 337 cursos.
Para esses, a moda foi o conceito 3, com 128 cursos.
Os demais cursos presenciais da região foram alocados nos conceitos 2 ( 87 cursos ), 4 ( 50 cursos ), 5 ( 16 cursos ) e 1 ( três cursos ). 28 cursos presenciais ficaram sem conceito ( SC ).
Na região Sudeste, 25 cursos foram oferecidos na modalidade a distância, sendo nenhum curso com conceito 1 ; 14 com conceito 2 ; sete com conceito 3 ; um com conceito 4 ; e nenhum com conceito 5.
Três cursos ficaram sem conceito ( SC ).
As instituições privadas concentraram 117 dos 144 cursos participantes da região Sul, 81,3 \% do total regional.
Desses, 53 cursos obtiveram conceito 3, o conceito modal.
Nessa combinação de categoria administrativa e grande região, 33 receberam o conceito 2 ; 18 receberam o conceito 4 ; um, o conceito 1 ; e nenhum obteve o conceito 5.
Doze cursos ficaram sem conceito ( SC ).
As instituições públicas na região Sul participaram com 27 cursos ( 18,8 \% dos cursos da região ), aos quais foram atribuídos os conceitos 2 ( dois cursos ), 3 ( oito cursos ), 4 ( 12 cursos ) e 5 ( cinco cursos ).
Nenhum curso da região Sul oferecido por IES públicas ficou com conceito 1 ou sem conceito ( SC ).
Na região Sul, 120 cursos foram oferecidos na modalidade presencial.
Para esses, a moda foi o conceito 3, com 51 cursos.
Os demais cursos presenciais da região foram alocados nos conceitos 4 ( 29 cursos ), 2 ( 25 cursos ), 5 ( cinco cursos ) e 1 ( nenhum curso ).
Ficaram sem conceito ( SC ) dez cursos.
Na modalidade a distância, foram oferecidos 24 cursos, sendo um curso com conceito 1, dez cursos com conceito 2 e dez com conceito 3 ( bimodal ), um curso com conceito 4 e nenhum curso com conceito 5.
Dois cursos ficaram sem conceito ( SC ) nessa região.
Na região Centro-Oeste, 25 dos 34 cursos participantes eram de IES privadas ( 73,5 \% em termos regionais ).
Desses, 11 cursos receberam conceito 2, o valor modal.
Dos demais cursos, sete receberam o conceito 3, um recebeu o conceito 4, um recebeu o conceito 1 e nenhum recebeu o conceito 5.
Cinco ficaram sem conceito ( SC ).
Quanto aos nove cursos oferecidos por instituições públicas na região Centro-Oeste ( 26,5 \% do total regional ), nenhum ficou com conceito 1 ; dois, com conceito 2 ; três, com conceito 3 ; três, com conceito 4 ; e um ficou com conceito 5.
Nenhum ficou sem conceito ( SC ) nessa região.
Também na região Centro-Oeste, a maioria dos cursos foi ofertada na modalidade presencial : 27 ( 79,4 \% ).
Para esses, a moda foram os conceitos 2 e 3, com 8 cursos cada ( 29,6 \% ) ( bimodal ).
Os demais cursos presenciais da região foram alocados nos conceitos 4 ( quatro cursos, 14,8 \% ), 5 ( um curso, 3,7 \% ) e 1 ( um curso, 3,7 \% ).
Cinco cursos presenciais da região ficaram sem conceito ( SC ).
Na modalidade a distância, foram oferecidos sete cursos, sendo cinco cursos com conceito 2 ( modal ), dois cursos com conceito 3, nenhum curso com os conceitos 1, 4 e 5.
Nenhum curso ficou sem conceito ( SC ) nessa região.
A título de síntese, no Gráfico 4.2, apresenta -se a distribuição do Conceito Enade, segundo a categoria administrativa da IES.
Nesse gráfico, como dito anteriormente, poligonais mais à esquerda representam Conceito Enade mais baixos, e poligonais mais à direita, Conceito Enade mais altos.
É importante observar que, para cada poligonal presente nesse gráfico, a porcentagem constante no eixo y refere -se ao universo de cursos oferecidos por instituições públicas ( poligonal azul ) ou privadas ( poligonal verde ), respectivamente.
Os cursos em IES públicas ( poligonal azul ) apresentam uma poligonal mais à direita do que os em IES privadas ( poligonal verde ) e, consequentemente, uma distribuição de Conceito Enade com valores maiores. 4.3 CONCEITOS POR ORGANIZAÇÃO ACADÊMICA E GRANDE REGIÃO Na Tabela 4.3, encontra -se a distribuição dos conceitos atribuídos aos cursos participantes do Enade 2023, na área de Engenharia de Produção, por organização acadêmica, segundo as grandes regiões brasileiras.
Como mostram os dados da Tabela 4.3, dos 638 cursos de Engenharia de Produção participantes, 247 eram oferecidos em universidades, 203 em centros universitários e 173 em faculdades ; 15 cursos foram oferecidos por CEFET/IF.
Essa distribuição corresponde a, respectivamente, 38,7 \%, 31,8 \%, 27,1 \% e 2,4 \% dos cursos.
De acordo com os dados apresentados na mesma tabela, de todos os 103 cursos avaliados com o conceito 4, 63 eram vinculados a universidades.
Os cursos vinculados a esse tipo de organização acadêmica tiveram o conceito 3 como conceito modal, atribuído a 97 cursos ( 39,3 \% ).
Os demais cursos ofertados por universidades receberam os conceitos 2 ( 45 cursos ), 5 ( 25 cursos ) e 1 ( dois cursos ). 15 cursos ficaram sem conceito ( SC ).
Dos cursos vinculados a centros universitários, o conceito modal também foi o 3, atribuído a 85 cursos ( 41,9 \% ). 18 cursos vinculados a esse tipo de organização acadêmica ficaram sem conceito ( SC ), e os outros cursos receberam os conceitos 2 ( 73 cursos ), 4 ( 23 cursos ), 5 ( dois cursos ) e 1 ( dois cursos ).
Dos 173 cursos mantidos por faculdades, 79 ( 45,7 \% ) receberam conceito 2, o conceito modal.
O segundo conceito mais frequentemente atribuído aos cursos oferecidos por essa categoria de IES foi o conceito 3, recebido por 50 cursos.
Na sequência, dez cursos obtiveram conceito 4 ; quatro cursos, o conceito 1 ; e nenhum curso, o conceito 5.
Dos cursos oferecidos em faculdades, 30 cursos ficaram sem conceito ( SC ).
Dos cursos oferecidos por CEFET/IF, o conceito foi bimodal, 3 e 4, atribuído a sete cursos ( 46,7 \% ) cada.
Nenhum curso vinculado a esse tipo de organização acadêmica ficou sem conceito ( SC ), ainda, um curso recebeu conceito 2 e nenhum curso obteve os conceitos 1 ou 5.
Considerando -se, separadamente, as regiões brasileiras, verifica -se que, na região Norte, as universidades participaram com oito cursos avaliados, aos quais foram atribuídos os seguintes conceitos : 1 a nenhum curso ; 2 a dois cursos ; 3, o conceito modal, a três cursos ; conceito 4 a dois cursos ; e 5 a um curso.
Nenhum curso ficou sem conceito ( SC ).
Os centros universitários da região Norte foram representados por nove cursos, dos quais nenhum curso obteve conceito 1 ; cinco cursos, conceito 2, o valor modal ; dois cursos, conceito 3 ; nenhum curso, o conceito 4 ; e um curso, o conceito 5.
Um curso ficou sem conceito ( SC ).
As faculdades participaram com dez cursos na região Norte, os quais receberam os conceitos 1 ( um curso ), 2 ( três cursos ), 3 ( um curso ), 4 ( um curso ) e 5 ( nenhum curso ).
Quatro cursos ficaram sem conceito ( SC ).
Não houve curso participante dos CEFET/IF nessa região.
Na região Nordeste, as universidades participaram com 34 dos 96 cursos da área de Engenharia de Produção oferecidos.
O conceito modal foi 3, atribuído a 14 cursos.
Outros dez cursos obtiveram o conceito 4 ; cinco cursos, o conceito 2 ; três cursos, o conceito 5 ; e nenhum curso, o conceito 1.
Dois cursos ficaram sem conceito ( SC ).
Os centros universitários contaram com 34 cursos participantes na região Nordeste, 14 dos quais receberam o conceito modal 2.
Para os demais cursos, foram atribuídos os conceitos 3 ( 11 cursos ), 4 ( cinco cursos ), nenhum curso ficou com conceito 5 ou 1.
Quatro cursos oferecidos por centros universitários no Nordeste ficaram sem conceito ( SC ).
As faculdades foram representadas por 27 cursos na região Nordeste, os quais receberam os conceitos 1 ( dois cursos ), 2 ( 20 cursos ), 3 ( um curso ), 4 ( nenhum curso ) e 5 ( nenhum curso ).
Quatro cursos ficaram sem conceito ( SC ).
Na região Nordeste, os CEFET/IF participaram com um curso, ao qual foi atribuído o conceito 3.
Na região Sudeste, as universidades concentraram 130 dos 337 cursos de Engenharia de Produção da região.
Dos cursos oferecidos em IES com este tipo de organização acadêmica, nessa região, o conceito modal foi 3, atribuído a 54 cursos.
Os demais cursos receberam os conceitos 4 ( 27 cursos ), 2 ( 23 cursos ), 5 ( 15 cursos ) e 1 ( dois cursos ).
Nove cursos ficaram sem conceito ( SC ).
Os centros universitários participaram com 99 cursos na região Sudeste, dos quais 41 obtiveram o conceito modal 3.
O conceito 2 foi atribuído a 35 cursos ; o conceito 4 a 12 cursos ; o conceito 5 a um curso ; e o conceito 1 a nenhum curso.
Dez cursos ficaram sem conceito ( SC ) nessa combinação de organização acadêmica e região.
Já as faculdades foram representadas por 97 cursos na região Sudeste, dos quais 12 ficaram sem conceito ( SC ).
O conceito modal 2 foi atribuído a 43 cursos.
Aos demais cursos foram atribuídos os conceitos 3 ( 34 cursos ), 4 ( sete cursos ), 1 ( um curso ) e 5 ( nenhum curso ).
Os CEFET/IF participaram com 11 cursos na região Sudeste, aos quais foram atribuídos os seguintes conceitos : 3 a seis cursos ; 4 a cinco cursos.
Nenhum curso ficou sem conceito ( SC ).
Dos 144 cursos da região Sul, 61 eram mantidos por universidades, para os quais o conceito modal foi 3, atribuído a 22 cursos.
Os demais cursos receberam os conceitos 4 ( 21 cursos ), 2 ( dez cursos ), 5 ( cinco cursos ) e 1 ( nenhum curso ).
Em IES com esse tipo de organização acadêmica, três cursos da região Sul ficaram sem conceito ( SC ).
Dos 54 cursos participantes de centros universitários da região Sul, três ficaram sem conceito ( SC ), e 28 cursos receberam conceito 3, o valor modal.
Os conceitos restantes ficaram assim distribuídos : 2 ( 16 cursos ), 4 ( seis cursos ), 1 ( um curso ) e 5 ( nenhum curso ).
Quanto aos 27 cursos vinculados a faculdades na região Sul, 11 cursos receberam o conceito 3, valor modal.
Os demais conceitos ficaram assim distribuídos : 2 ( nove cursos ), 4 ( um curso ), 5 ( nenhum curso ) e 1 ( nenhum curso ).
Dos 27 cursos, seis ficaram sem conceito ( SC ).
Os CEFET/IF participaram com dois cursos na região Sul, aos quais foram atribuídos o conceito 4.
Na região Centro-Oeste, 14 dos 34 cursos oferecidos eram mantidos por universidades.
Para esse tipo de organização acadêmica, o conceito modal foi 2, atribuído a cinco cursos.
Em relação às universidades, um curso da região ficou sem conceito ( SC ).
Os demais cursos receberam os conceitos 3 ( quatro cursos ), 4 ( três cursos ), 5 ( um curso ) e 1 ( nenhum curso ).
Os centros universitários da região Centro-Oeste contaram com sete cursos : três cursos foram avaliados com os conceitos 3 e 2 ( bimodal ) ; um curso, com o conceito 1.
Nenhum curso ficou sem conceito ( SC ).
Dos 12 cursos oferecidos em faculdades na região Centro-Oeste, quatro cursos receberam o conceito 2 ( modal ).
Os conceitos restantes ficaram assim distribuídos : 3 ( três cursos ), 4 ( um curso ), 5 e 1 ( nenhum curso ).
Nessa região, quatro cursos oferecidos em faculdades ficaram sem conceito ( SC ).
No que se refere aos CEFET/IF da região Centro-Oeste, um curso foi avaliado, alcançando o conceito 2.
De maneira semelhante à forma como se procedeu em relação à seção anterior, esta seção \textbf{encerra} -se com o Gráfico 4.3, que apresenta a distribuição do Conceito Enade dos cursos de Engenharia de Produção segundo a organização acadêmica da IES.
No Gráfico 4.3, apresenta -se a distribuição do Conceito Enade segundo a organização acadêmica da IES.
A leitura desse último gráfico é semelhante aos anteriores, seguindo -se a lógica de que poligonais deslocadas para a esquerda correspondem aos conceitos mais baixos ( 1, 2 e 3 ) e poligonais deslocadas para a direita correspondem a Conceitos Enade mais altos.
Nas universidades ( linha azul ), todos os cursos obtiveram conceito 5 ou abaixo disso, já que a poligonal cumulativa alcança 100 \% para o conceito 5.
Além disso, os cursos de CEFET/IF ( linha amarela ) apresentam uma poligonal mais à direita do que as demais, denotando uma melhor distribuição de Conceitos Enade.
No outro extremo, encontram -se os cursos mantidos por faculdades ( linha vermelha ), apresentando as respectivas poligonais mais à esquerda e, consequentemente, uma distribuição de Conceitos Enade com os menores valores.
Os cursos mantidos por universidades ( linha azul ) e CEFET/IF ( linha amarela ) apresentam poligonais mais centrais e, consequentemente, uma distribuição de Conceitos Enade com os valores intermediários.

% matched lemmas: ajudar, encerrar, laranja, traçar
\end{textsample}
